\section{Markov transfer compatible with the entire inverse limit}
\label{sec:cpe-transferencia-limite-inverso}

Let \(\rho_m:H_m\twoheadrightarrow Y_m\) be the observable quotient at depth
\(m\), and let \(L_m(y,y')\) be the transition kernel already constructed on
\(Y_m\). For each finite fibre \(\rho_m^{-1}(y)\), let
\(\nu_{m,y}\) be the conditional uniform measure. The nonadic regularity of
the extensions allows these measures to be chosen so that their pushforward under
truncation is \(\nu_{m-1,\bar y}\).

Define, for \(h\in H_m\),
\begin{equation}
 \mathcal K_m(h,A)
 :=\sum_{y'\in Y_m}
 L_m(\rho_m(h),y')\,
 \nu_{m,y'}\bigl(A\cap\rho_m^{-1}(y')\bigr).
 \label{eq:cpe-kernel-lift-uniforme}
\end{equation}

\begin{theorem}[Strongly additive elevation and Markov limit]
\label{thm:cpe-kernel-proyectivo}
The kernels \(\mathcal K_m\) satisfy
\begin{equation}
 \sum_{\rho_m(g)=y'}\mathcal K_m(h,\{g\})
 =L_m(\rho_m(h),y'),
 \label{eq:cpe-agregabilidad-fuerte}
\end{equation}
independently of the representative \(h\) of the fibre. Moreover, they commute with
the truncations. Therefore, there exists a unique Markov kernel
\(\mathcal K_\infty\) on the inverse limit whose level-\(m\) marginal
is \(\mathcal K_m\). Its projection by \(\rho_\infty\) is the observable
kernel \(L_\infty\).
\end{theorem}

\begin{proof}
The sum of the left-hand side of
\eqref{eq:cpe-agregabilidad-fuerte} is
\[
 L_m(\rho_m(h),y')
 \sum_{g\in\rho_m^{-1}(y')}\nu_{m,y'}(\{g\})
 =L_m(\rho_m(h),y').
\]
Compatibility of \(L_m\), \(\rho_m\), and the conditional measures proves the truncation square. The cylindrical distributions obtained by iterating \(\mathcal K_m\) are consistent; the extension theorem for projective systems of finite spaces produces \(\mathcal K_\infty\), and uniqueness follows because the cylinders generate the sigma-algebra of the limit.
\end{proof}

The construction preserves microscopic stillness and complete labels
in the conditioned measure. The visible average does not replace the
enriched state and it is not identified with the Topological Phase Kernel.

